
\subsection{Creation of colorboxes}
\newtcolorbox{rb}{%
  width=0.7\linewidth,
  center,
  halign=flush center,
  colback=resultsbackground,
  colframe=resultscolor,
  top=10mm,    % Innenabstand oben verringern (Standard ist ~2mm)
  bottom=1mm  % Innenabstand unten verringern
}
\verb|top,bottom| sind Parameter für den Abstand von Text und Rand. \verb|afterskip,beforeskip| für den Abstand der Box vom Fließtext.
\begin{rb}
  \verb|top=10mm,bottom=1mm|
\end{rb}

\subsection{Mathemodus innerhalb tcb -- sinnlos}

\newtcolorbox{testbox}{width=0.7\linewidth,center,halign=center,colback=white,colframe=black}
\newtcolorbox{testboxlong}{width=\linewidth,center,halign=center,colback=white,colframe=black}
Testboxen sind grauweiß und mit \verb|center,halign=center| ausgestattet. Ersteres centered die Box auf der Seite, zweites ist für Textausrichtung innerhalb der Box.

Displaystyle mit \verb|\[\]|:
\begin{testbox}
  \[E_{3p}=\qty{-3.0256(28)}{\electronvolt}\frac12\]
\end{testbox}
Ergo: Zwar gutes kerning, aber vertikale Zentrierung funktioniert nicht. 
\begin{testbox}
  DIes ist ein Test
\end{testbox}

nur inline \verb|$$|:
\begin{testbox}
  $E_{3p}=\qty{-3.0256(28)}{\electronvolt}\frac12$
\end{testbox}
Zwar vertikal zentriert, aber schlechtes kerning.

Inline mit \verb|$$| und \verb|\displaystyle|:
\begin{testbox}
  $\displaystyle
  E_{3p}=\qty{-3.0256(28)}{\electronvolt}\frac12$
\end{testbox}
Es hat sich nichts geändert, weiterhin schlechtes kerning, dafür sind Brüche groß.

mit \verb|aligned| und \verb|$$|:
\begin{testbox}
  $\begin{aligned}
  E_{3p}=\qty{-3.0256(28)}{\electronvolt}\frac12
    \end{aligned}$
\end{testbox}
gutes Kerning, und vertikal zentriert. 

\verb|aligned| mit \verb|\[\]|:
\begin{testbox}
  \[\begin{aligned}
  E_{3p}=\qty{-3.0256(28)}{\electronvolt}\frac12
    \end{aligned}\]
\end{testbox}
es scheißt wieder vertikal rein. 

Nun mal mit \verb|ams algin*|:
\begin{tcolorbox}[ams align*]
    E_{3p}=\qty{-3.0256(28)}{\electronvolt}\frac12
\end{tcolorbox}

default tcolorbox mit alignment und Displaystyle mit \verb|\[\]|:
\begin{tcolorbox}[width=0.7\textwidth, center, valign=center]
    \[E_{3p}=\qty{-3.0256(28)}{\electronvolt}\]
\end{tcolorbox}

normal math mode \verb|\(\)|
\begin{center}
$E_{3p}=\qty{-3.0256(28)}{\electronvolt}$
\end{center}
display math mode \verb|\[\]|
\[E_{3p}=\qty{-3.0256(28)}{\electronvolt}\]

\subsection{math Optionen für tcb}
\begin{tcolorbox}[width=0.7\textwidth,ams align*,colback=yellow!10!white,colframe=red!50!black]
\sum\limits_{n=1}^{\infty} \frac{1}{n} &= \infty. \quad \sum_{n=1}^{\infty}\frac1n=\infty\qquad a=b
\end{tcolorbox}

ams algin*
\begin{tcolorbox}[width=0.7\textwidth,ams align*,colback=yellow!10!white,colframe=red!50!black]
a=b
\end{tcolorbox}

\verb|math,halign=center|
\begin{tcolorbox}[width=0.7\textwidth,math,halign=center,colback=yellow!10!white,colframe=red!50!black]
a=b
\end{tcolorbox}

\verb|math|
\begin{tcolorbox}[width=0.7\textwidth,math,colback=yellow!10!white,colframe=red!50!black]
a=b
\end{tcolorbox}
YEAAAH, macht eigentlich genau das was ich will.
\verb|math,halign=flush center|
\begin{tcolorbox}[width=0.7\textwidth,math,halign=flush center,colback=yellow!10!white,colframe=red!50!black]
a=b
\end{tcolorbox}

\verb|math,center,halign=flush center|
\begin{tcolorbox}[width=0.7\textwidth,math,center,halign=flush center,colback=yellow!10!white,colframe=red!50!black]
E_{3p}=\qty{-3.0256(28)}{\electronvolt}\frac12
\end{tcolorbox}
ICH WEINE. WARUM HABE ICH NICHT GENAUER GELESEN. FLUSH CENTER WAR DAS RICHTIGE.

\verb|ams align* with &&|
\begin{tcolorbox}[ams align*]
    A&&B&&C
\end{tcolorbox}

\subsection{tcbhighmath}
normal
\[\tcbhighmath{E_{3p}=\qty{-3.0256(28)}{\electronvolt}\frac12}\]

in align
\begin{align}
\tcbhighmath{\sum\limits_{n=1}^{\infty} \frac{1}{n}} &= \infty.\\
\int x^2 ~\text{d}x &= \frac13 x^3 + c.
\end{align}

falign design aus nem stackexchange post
\begin{flalign}
  \begin{gathered}
    A  \vphantom{\sum \limits _{n=0} ^{N-1}}\\
    B \vphantom{\sum \limits _{n=0} ^{N-1}}
  \end{gathered}
  &&
  \tcbhighmath{\begin{aligned}\label{equationAB}
    A &= \sum \limits _{n=0} ^{N-1} a[n]\\
    B &= \sum \limits _{n=0} ^{N-1} b[n]
  \end{aligned}}
  &&
\end{flalign}

\verb|borderline={<linewidth>}{<Versatz zu bounding box>}{<color>}|
\tcbset{highlight math style={enhanced,colback={resultsbackground},borderline={1pt}{-4mm}{resultscolor},left=0pt,right=0pt,top=0pt,bottom=0pt,extrude left by=0pt,extrude right by=0pt,extrude top by=0pt,extrude bottom by=0pt,boxsep=\fboxsep}}

Inline Ergebnisse: $\tcbhighmath{p^\text{Luft}=\qty{1020.0(0.6)}{\hecto\pascal}$, $p^\text{Argon}=\qty{1020.1(0.6)}{\hecto\pascal}.}$

\subsubsection{Results Design}

\tcbset{highlight math style={enhanced, colback={resultsbackground},colframe={resultscolor},left=0pt,right=0pt,top=0pt,bottom=0pt,extrude left by=1pt,extrude right by=1pt,extrude top by=1pt,extrude bottom by=1pt,boxsep=\fboxsep}}

Inline Ergebnisse: $\tcbhighmath{p^\text{Luft}=\qty{1020.0(0.6)}{\hecto\pascal}$, $p^\text{Argon}=\qty{1020.1(0.6)}{\hecto\pascal}.}$

Oder in einer \verb|align &&|-Umgebung für Formel-Fehler-Ergebnis:
\begin{align}
  F=ma&&\Delta F=F\sqrt{\left(\frac{\Delta m}{m}\right)^2+\left(\frac{\Delta a}{a}\right)^2} && \tcbhighmath{F=\qty{1}{\newton}}
\end{align}
Optional \verb|[boxsep=2mm]| oder so bei Ergebnissen in \verb|align| einstellen, sodass die Box größer ist als innerhalb des Textes.
\begin{align}
  F=ma&&\Delta F=F\sqrt{\left(\frac{\Delta m}{m}\right)^2+\left(\frac{\Delta a}{a}\right)^2} && \tcbhighmath[boxsep=2mm]{F=\qty{1}{\newton}}
\end{align}

\paragraph{Preresults}
Optional \verb|[colback=white]| einstellen.
\begin{align}
  F=ma&&\Delta F=F\sqrt{\left(\frac{\Delta m}{m}\right)^2+\left(\frac{\Delta a}{a}\right)^2} && \tcbhighmath[colback=white,boxsep=2mm]{F=\qty{1}{\newton}}
\end{align}

\subsubsection{results alignment}
align mit \verb|&&&&| vor letzter Spalte
\begin{align}
  F=ma&&\Delta F=F\sqrt{\left(\frac{\Delta m}{m}\right)^2+\left(\frac{\Delta a}{a}\right)^2} &&&& \tcbhighmath[colback=white,boxsep=2mm]{F=\qty{1}{\newton}}
\end{align}

flalign ohne startendes \verb|&&|, damit Ergebnisse immer rechtsbündig sind. 
\begin{flalign}
  F=ma&&\Delta F=F\sqrt{\left(\frac{\Delta m}{m}\right)^2+\left(\frac{\Delta a}{a}\right)^2} && \tcbhighmath[colback=white,boxsep=2mm]{F=\qty{1}{\newton}}
\end{flalign}

\verb|&&|
\begin{flalign}
  &&F=ma&&\Delta F=F\sqrt{\left(\frac{\Delta m}{m}\right)^2+\left(\frac{\Delta a}{a}\right)^2} && \tcbhighmath[colback=white,boxsep=2mm]{F=\qty{1}{\newton}}
\end{flalign}

Ausrichtung der tcb mit \verb|&\tcbhighmath| und \verb|&&&&| in der zweiten Zeile 
\begin{align}
    B &= \frac{U_\text{ind, max}}{A_0 n_f 2\pi f_D} &
    \Delta B &= B\sqrt{\left(\frac{\Delta U_\text{ind, max}}{U_\text{ind, max}}\right)^2+\left(\frac{\Delta f_D}{f_D}\right)^2} &&\tcbhighmath{B=\qty{52(3)}{\micro\tesla}} \\
    &&&& &\tcbhighmath{B_\parallel=\qty{53.4(0.4)}{\micro\tesla}}
\end{align}
